\section{Motivation}
\label{sec:motivation}

Based on the workload characterization above, we identify two challenges in
deploying real-time depth completion on a heterogeneous edge platform.  Each
challenge exposes an opportunity that directly motivates one part of
\system's design.

\paragraph{Challenge 1: Monolithic execution misses the real-time budget.}

The Transformer encoder and CNN decoder have different operator mixes and
hardware preferences, yet a conventional deployment executes both stages
sequentially on one processor.  The resulting end-to-end latency includes the
full encoder, decoder, and runtime path for every frame.  Merely moving the
decoder to another device does not remove this serialization if communication
and accelerator execution remain on the critical path.

\paragraph{Opportunity 1: Task-level heterogeneous pipelining.}
This opportunity follows from the clear stage separation.
The clear encoder--decoder boundary permits a coarse-grained partition: the
Thor executes the encoder while the FPGA executes the decoder.  Across a
stream of frames, queued communication and double buffering can overlap the
encoder of frame $i+1$ with transfer and decoding of frame $i$.  This turns
device specialization into steady-state pipeline throughput rather than an
additional synchronous offload cost.

\paragraph{Challenge 2: End-to-end low precision can degrade metric depth.}

Low precision is beneficial on both sides of the partition.  In particular,
selected MLP computation in the Thor-side encoder and the FPGA-side decoder
can be executed at reduced precision.  A uniform quantization policy, however,
must simultaneously accommodate Transformer activation distributions,
depth-regression sensitivity, and the numerical contract at the cross-device
boundary.  Aggressive quantization can therefore save computation and
bandwidth while introducing scale errors or local depth artifacts.

\paragraph{Opportunity 2: Coordinated depth- and hardware-aware precision.}
This opportunity addresses the tension between efficiency and depth quality.
Rather than treating encoder quantization, feature transfer, and decoder
quantization as independent optimizations, \system{} coordinates their
precision choices.  Depth-aware quantization (DAQ) targets the layers and
signals that most affect metric-depth accuracy, while hardware-aware precision
choices map the encoder MLP and decoder datapath to the efficient modes of
Thor and FPGA, respectively.  The same boundary representation also reduces
the cost of cross-device feature transfer.

\begin{figure}[t]
  \centering
  \figslot{figures/motivation-evidence.pdf}{1.40in}{%
    Motivation evidence\\[5pt]
    (a) Sequential versus asynchronous frame schedule\\
    (b) Accuracy sensitivity of encoder-MLP and decoder quantization\\
    Use measured latency and depth error only}
  \caption{Evidence motivating \system. (a) A synchronous offload leaves
  communication and decoder execution on the critical path, whereas a
  task-level schedule overlaps work across frames. (b) Low-precision choices
  in both the encoder MLP and decoder must be coordinated to avoid excessive
  metric-depth error.}
  \Description{Two panels comparing sequential and asynchronous execution and
  showing measured depth-error sensitivity to low precision in the encoder
  MLP and decoder.}
  \label{fig:motivation}
\end{figure}
